% Generated draft: requires mathematical and source review.
\documentclass[11pt]{article}
\usepackage[margin=25mm]{geometry}
\usepackage{fontspec}
\setmainfont{Latin Modern Roman}
\usepackage{amsmath,amssymb,mathtools,unicode-math}
\setmathfont{Latin Modern Math}
\usepackage{enumitem,needspace,xurl,hyperref}
\setlength{\parindent}{0pt}
\setlength{\parskip}{5pt}
\setlength{\emergencystretch}{3em}
\newcommand{\N}{\mathbb{N}}
\newcommand{\Z}{\mathbb{Z}}
\newcommand{\Q}{\mathbb{Q}}
\newcommand{\R}{\mathbb{R}}
\newcommand{\C}{\mathbb{C}}
\newcommand{\F}{\mathbb{F}}
\newcommand{\A}{\mathbb{A}}
\newcommand{\PP}{\mathbb{P}}
\newcommand{\E}{\mathbb{E}}
\newcommand{\eps}{\varepsilon}
\newcommand{\dd}{\,\mathrm{d}}
\begin{document}
\section*{UnsolvedMath ID: 30004736}
\textbf{Source problem code:} catalogue entry
\subsection{Definitions and mathematical statement}
The supplied question concerns a real-analytic convex billiard domain \(\Omega\), its perimeter \(P=|\partial\Omega|\), and its length spectrum. The record does not explicitly specify the ambient space or the convention for the length spectrum. The following gives a precise conditional formulation under the usual planar Euclidean interpretation; these conventions require confirmation against the source.

Assume that \(\Omega\subset\R^2\) is a bounded open convex domain with regular real-analytic boundary. A periodic billiard trajectory is a closed polygonal path with finitely many nonzero edges, vertices on \(\partial\Omega\), and edge interiors in \(\Omega\), satisfying specular reflection at each vertex. Explicitly, if \(v_-\) and \(v_+\) are the incoming and outgoing unit velocity vectors and \(n\) is the outward unit normal, then
\[
 v_+=v_- -2\langle v_-,n\rangle n.
\]
Under the convention that repeated traversals are included and that lengths are recorded as a set, define
\[
\mathcal L(\Omega)=\{\ell(\gamma):\gamma\text{ is a periodic billiard trajectory in }\Omega\},
\]
where \(\ell(\gamma)\) is the total Euclidean length of the polygonal path. Define its set of finite accumulation points by
\[
\operatorname{Acc}\mathcal L(\Omega)
=\left\{a\in\R:\ \forall\varepsilon>0,\quad
\mathcal L(\Omega)\cap\bigl((a-\varepsilon,a+\varepsilon)\setminus\{a\}\bigr)\ne\varnothing\right\}.
\]
Thus accumulation concerns distinct length values, rather than repeated occurrences of the same value.

\textbf{Question under these conventions.} For every such \(\Omega\), does
\[
\operatorname{Acc}\mathcal L(\Omega)
\subseteq \{mP:m\in\Z,\ m\ge 1\}
\]
hold?

The wording ``the only accumulation points'' could additionally intend the stronger assertion
\[
\operatorname{Acc}\mathcal L(\Omega)
=\{mP:m\in\Z,\ m\ge 1\}.
\]
The supplied record does not distinguish these readings. It also does not explicitly require strict convexity or everywhere positive curvature, specify whether only primitive trajectories are counted, or address degenerate or grazing trajectories. These qualifications remain unresolved; no stronger convexity hypothesis is imposed here.
\subsection{Short English statement}
Under an explicitly provisional planar billiard convention, ask whether every finite accumulation point of the set of periodic trajectory lengths is a positive integer multiple of the boundary perimeter. The source leaves some geometric and spectral conventions unspecified, and its wording may intend equality with all positive perimeter multiples.
\subsection{Sources}
\begin{itemize}
\item Local UnsolvedMath registry, supplied complete source record: permanent ID \(30004736\), title \emph{Accumulation Points of Convex Billiard Length Spectra}, exact source problem code \texttt{catalogue entry}. Collection: \emph{Oberwolfach Reports Open Problems}. No individual proposer is supplied; the proposed year is recorded as 2021.
\item Supplied report citation: \emph{Dynamische Systeme (2022)}, Oberwolfach Reports, pp.~1735--1803, DOI \texttt{10.4171/owr/2021/33}. \url{https://doi.org/10.4171/owr/2021/33}. The supplied citation says 2022, whereas the DOI and proposed-year metadata contain 2021; both are retained without resolving the dating.
\item Supplied literature source, labeled ``formal Birkhoff-normal-form rigidity'': \url{https://msp.org/apde/2012/5-5/apde-v5-n5-p06-s.pdf}.
\item Supplied literature source, labeled ``analytic symmetric rigidity'': \url{https://annals.math.princeton.edu/2009/170-1/p06}.
\item Supplied literature source, labeled ``almost-every ellipse uniqueness'': \url{https://doi.org/10.1007/s00222-025-01397-y}.
\item Literature assessment checked at \texttt{2026-08-22}: this is supplied dated curation, not independently verified. The record labels the problem \texttt{partially\_solved} and states that analytic symmetry classes have spectral rigidity and formal nonisometric Birkhoff-normal-form germs are known, while general analytic realization and convexity-detection questions remain open. This assessment is preserved as supplied; its relationship to the present accumulation-point question is not established in the record. The background expressly requests independent verification of time-sensitive status and cited claims.
\item The supplied statement status is \texttt{corrected\_verified}, with a claim of manual semantic, typographic, and LaTeX review. This is source metadata and does not independently establish mathematical correctness or resolve the formulation issues above.
\end{itemize}
\end{document}
